\section{Related Work}
\label{sec:related-work}

\subsection{Tool-Using Language Agents and Agentic Reinforcement Learning}

ReAct demonstrated that interleaving language reasoning and actions can improve
interactive decision making \citep{yao2023react}, while ToolLLM scaled
instruction tuning to a large collection of real-world APIs
\citep{qin2024toolllm}. Stateful benchmarks subsequently made it possible to
evaluate whether an agent completes a goal rather than merely predicts a
function name. $\tau$-bench evaluates tool-agent-user interaction in realistic
domains, and $\tau^2$-Bench extends the setting to dual-control conversations in
which both the user and the agent affect the trajectory
\citep{yao2024taubench,barres2025tau2bench}.

Reinforcement-learning approaches target limitations of supervised tool traces.
ToolRL studies reward design for tool learning \citep{qian2025toolrl};
AgentGym-RL trains agents over long-horizon, multi-turn interactions
\citep{xi2025agentgymrl}; and VerlTool and RLFactory provide infrastructure for
multi-turn rollout and post-training \citep{jiang2025verltool,chai2025rlfactory}.
Many recent systems use group-relative policy optimization, popularized in
open-model reasoning by DeepSeekMath \citep{shao2024deepseekmath}. Our work
shares the executable, multi-turn training setting but studies whether recovery
mechanisms remain identifiable after matching rollout estimators and interaction
budgets.

\subsection{Robustness, Noise, and Open-World Tool Use}

Robustness studies show that function-call performance is sensitive to
natural-language reformulations and changes in the available tool set
\citep{rabinovich2025robustness}. AgentNoiseBench broadens this analysis to
user-side and tool-side noise \citep{wang2026agentnoisebench}. ToolBench-X
organizes recoverable tool-environment unreliability into specification drift,
invocation error, execution failure, output drift, and cross-source conflict,
with executable multi-step tasks \citep{tian2026unreliability}. OpenAgent
studies query, action, observation, and domain shifts in an open-world hierarchy
and finds that static training remains brittle under environmental change
\citep{lv2026openagent}. Sim-to-real work further cautions that simulator success
need not transfer to realistic tool execution \citep{zhou2026simulationlies}.

NoisyAgent introduces a noise curriculum for perturbed environments
\citep{chen2026noisyagent}, and EnvFactory studies executable environment
synthesis and robust RL at scale \citep{xu2026envfactory}. These methods show
that exposure to diverse environments can improve robustness. Our result is a
boundary condition: in a strong-policy regime, environment difficulty can be
dominated by the distance between the required recovery path and the turn
budget. Floor and ceiling effects can then hide or mimic algorithmic gains. We
therefore pre-calibrate the baseline into a fixed acceptance band before
freezing the discriminative suite.

\subsection{Failure Recovery and Reliability Evaluation}

Fission-GRPO, accepted to ACL 2026, directly trains recovery by converting
execution errors into on-policy corrective instances
\citep{zhang2026fissiongrpo}. It augments failed trajectories with diagnostic
feedback from a fine-tuned error simulator and resamples recovery rollouts,
reporting gains on BFCL v4 Multi-Turn, $\tau$-bench, and $\tau^2$-Bench. Our
failure-prefix branching is motivated by the same on-policy principle but uses
typed executable failures and state verifiers, applies an explicit recovery
budget, and adds fault-group minimax weighting. Crucially, our estimator-matched
ablation does not establish an advantage for branching: the MRPO versus
no-fission comparison changes sign across seeds and its pooled paired interval
contains zero. This is not an official reproduction of Fission-GRPO; it shows
that the general branching mechanism does not automatically transfer to our
controlled fault families and training regime.

Recent benchmarks make recovery itself the evaluation target. ToolMaze uses
DAG-based task complexity and a two-by-two taxonomy of explicit versus implicit
and transient versus permanent perturbations; it reports that implicit semantic
failures reduce perturbation recovery rate by roughly 37\%\footnote{ToolMaze:
\url{https://arxiv.org/abs/2606.05806}.}. ReliabilityBench evaluates a reliability
surface over repeated-execution consistency, semantic perturbation intensity,
and controlled fault intensity, including timeouts, rate limits, partial
responses, and schema drift\footnote{ReliabilityBench:
\url{https://arxiv.org/abs/2601.06112}.}. These evaluation-oriented works align
with our separation of unconditional task success, fault exposure, conditional
recovery, and recovery cost. Our additional focus is training-time
identifiability: equalizing the number of primary rollouts, the advantage
grouping unit, and the recovery budget before attributing a behavioral change
to a recovery mechanism.

Recoverability can also be summarized through regularity or risk measures. The
ERR measure studies whether tool-augmented agents recover from failures
\citep{vuddanti2026recoverability}, while ToolChain-CRC controls trajectory risk
under retrieval and tool-use drift \citep{opoku2026toolchaincrc}. In our
protocol, recovery requires verifier-confirmed task success and at least one
non-faulted action after the final observed failure, avoiding credit for a
hidden successful commit followed by immediate termination.

\subsection{Rollout Allocation and Worst-Group Optimization}

Tree-based agent training can spend most of its budget on redundant branches.
InfoTree casts rollout selection under a fixed budget as an informativeness
problem \citep{hu2026infotree}. \method{} instead scores failure prefixes using
estimated frequency, non-recovery probability, uncertainty, group weakness, and
expected cost. Our experiments show why generated-token accounting alone is
insufficient: primary rollout count and advantage grouping must also be matched
to identify a mechanism effect.

Worst-group objectives are closely related to group distributionally robust
optimization, which upweights high-loss groups to improve worst-case
generalization \citep{sagawa2020groupdro}. We use fault families as groups and
retain a clean-group floor. Group DRO itself is not a contribution. Moreover,
the three-seed minimax ablation is consistently unfavorable: removing minimax
weighting improves the primary-suite result, with a pooled paired interval below
zero for MRPO minus no-minimax. We therefore present minimax weighting as a
tested mechanism with a negative result, not as an empirically validated
improvement.

\paragraph{Positioning.}
The nearest work contributes noisy-environment training, recovery-instance
construction, and reliability evaluation. Our contribution is an auditable
protocol for determining whether such mechanisms cause behavioral improvement:
pre-calibrated and frozen executable tasks, estimator-matched ablations,
multi-seed task-paired inference, and clean/fault regression gates. The protocol
finds no reproducible MRPO advantage in the studied regime and identifies a
stable negative effect from minimax weighting. This narrows the conditions under
which structured recovery objectives should be expected to help, rather than
claiming a new state of the art.
